% Compile with xelatex or tectonic: SL_third_order_K1_proof.tex
\documentclass[11pt]{article}
\usepackage[margin=1in]{geometry}
\usepackage{amsmath,amssymb}
\newtheorem{theorem}{Theorem}
\newtheorem{lemma}[theorem]{Lemma}
\newtheorem{corollary}[theorem]{Corollary}
\newtheorem{definition}[theorem]{Definition}
\newtheorem{remark}[theorem]{Remark}
\newcommand{\R}{\mathbb{R}}

\title{A Strict Proof of the Even Minimal-Solution Anchor $K(1)=e/4$}
\author{BVE research project}
\date{2026-08-25}

\begin{document}
\maketitle

\begin{abstract}
The repository previously recorded the even minimal-solution constant
\(K(1)=e/4\) as a high-precision numerical conjecture.  This note proves the
claim directly from the defining backward recurrence.  An exact factorial
scaling reduces the third-order recurrence to a first-order recurrence for a
second difference.  The finite backward solution, its fixed-index limit, the
unique minimal branch, and the asymptotic tail are all obtained by positive
factorial series.  No numerical fit, formal asymptotic ansatz, or unchecked
computer algebra is used.
\end{abstract}

\section{Statement}

For \(j\geq 3\), set
\begin{align*}
P_j&=8j^2-4j+\frac{j}{j-1},\\
Q_j&=4j(j-1)(2j-1)(2j-3)+4j(2j-3),\\
R_j&=4j(j-2)(2j-3)(2j-5).
\end{align*}
Consider
\begin{equation}
\mu_j=P_j\mu_{j-1}-Q_j\mu_{j-2}+R_j\mu_{j-3}.
\tag{R}
\end{equation}
For \(N\geq 3\), define the backward solution by
\[
\mu_{N+1}^{(N)}=1,\qquad \mu_N^{(N)}=\mu_{N-1}^{(N)}=0,
\]
and use (R) backwards for \(j=N+1,N,\ldots,3\).  Whenever the denominator is
nonzero, put \(\widehat\mu_k^{(N)}=\mu_k^{(N)}/\mu_0^{(N)}\).

\begin{theorem}[Strict \(K(1)\) anchor]
For every fixed \(k\geq 0\), the limit
\(\mu_k^*=\lim_{N\to\infty}\widehat\mu_k^{(N)}\) exists.  It is the unique
normalized minimal solution and has the exact representation
\begin{equation}
\boxed{\displaystyle
\mu_k^*=2e(2k)!\sum_{r=k+2}^{\infty}
\frac{r-k-1}{(2r-1)!}}.
\tag{1}
\end{equation}
Moreover,
\begin{equation}
\boxed{\displaystyle \lim_{j\to\infty}j^3\mu_j^*=\frac{e}{4}}.
\tag{2}
\end{equation}
\end{theorem}

\section{Exact factorization}

Define
\[
v_j=\frac{\mu_j}{(2j)!},
\qquad
c_j=\frac{1}{2(j-1)(2j-1)}.
\]
Direct cancellation gives
\[
\frac{P_j(2j-2)!}{(2j)!}=2+c_j,\qquad
\frac{Q_j(2j-4)!}{(2j)!}=1+2c_j,\qquad
\frac{R_j(2j-6)!}{(2j)!}=c_j.
\]
Thus (R) is equivalent to
\begin{equation}
v_j=(2+c_j)v_{j-1}-(1+2c_j)v_{j-2}+c_jv_{j-3}.
\tag{3}
\end{equation}
Set
\[
d_j=v_j-2v_{j-1}+v_{j-2}.
\]
Equation (3) factors as
\begin{equation}
d_j=c_jd_{j-1}.
\tag{4}
\end{equation}
Since
\[
\frac{j/(2j)!}{(j-1)/(2j-2)!}
=\frac{1}{2(j-1)(2j-1)}=c_j,
\]
all solutions of (4) have the form
\begin{equation}
d_j=C\frac{j}{(2j)!}
\tag{5}
\end{equation}
for a constant \(C\).

\section{Finite backward solutions and the fixed-index limit}

For the terminal solution, \(v_N^{(N)}=v_{N-1}^{(N)}=0\) and
\(v_{N+1}^{(N)}=1/(2N+2)!\).  Hence
\[
d_N^{(N)}=0,\qquad d_{N+1}^{(N)}=\frac{1}{(2N+2)!}.
\]
Backward use of (4) gives
\begin{equation}
d_j^{(N)}=\frac{j}{(N+1)(2j)!},
\qquad 2\leq j\leq N+1.
\tag{6}
\end{equation}
Twice summing the second differences in (6), with the two terminal zeros,
gives the exact finite formula
\begin{equation}
\mu_j^{(N)}
=\frac{(2j)!}{2N+2}\sum_{r=j+2}^{N}
\frac{r-j-1}{(2r-1)!},
\qquad 0\leq j\leq N-1.
\tag{7}
\end{equation}
In particular,
\[
\mu_0^{(N)}=\frac{1}{2N+2}
\sum_{r=2}^{N}\frac{r-1}{(2r-1)!}>0
\]
for every \(N\geq 3\).  Therefore the prescribed normalization is always
valid, not merely eventually.

For fixed \(k\), (7) yields
\[
\widehat\mu_k^{(N)}=(2k)!
\frac{\displaystyle\sum_{r=k+2}^{N}(r-k-1)/(2r-1)!}
{\displaystyle\sum_{r=2}^{N}(r-1)/(2r-1)!}.
\]
Both sums are positive and convergent.  Moreover,
\begin{align*}
\sum_{r=2}^{\infty}\frac{r-1}{(2r-1)!}
&=\frac12\sum_{r=2}^{\infty}
\left(\frac{1}{(2r-2)!}-\frac{1}{(2r-1)!}\right)\\
&=\frac12\bigl((\cosh 1-1)-(\sinh 1-1)\bigr)
=\frac{1}{2e}.
\end{align*}
Taking the limit proves (1) and \(\mu_0^*=1\).

\section{Minimality and uniqueness}

Define
\[
\phi_j=\sum_{r=j+2}^{\infty}
\frac{r(r-j-1)}{(2r)!}.
\]
The series is absolutely convergent and tends to zero.  Direct telescoping of
the tails gives
\[
\phi_j-2\phi_{j-1}+\phi_{j-2}=\frac{j}{(2j)!}.
\]
Consequently every solution of (3) has the form
\begin{equation}
\frac{\mu_j}{(2j)!}=A+Bj+C\phi_j.
\tag{8}
\end{equation}
The minimal branch is the subfactorial branch, equivalently the branch for
which \(\mu_j/(2j)!\to 0\).  By (8), this forces \(A=B=0\), so the minimal
solution space is one-dimensional.  Since
\[
\phi_0=\sum_{r=2}^{\infty}\frac{r(r-1)}{(2r)!}=\frac{1}{4e},
\]
the condition \(\mu_0=1\) fixes \(C=4e\).  The resulting formula is exactly
(1), so the backward limit is the unique normalized minimal solution.

\section{Asymptotic constant}

Write \(r=j+\ell+1\) in (1):
\begin{equation}
\mu_j^*=2e\sum_{\ell=1}^{\infty}
\ell\frac{(2j)!}{(2j+2\ell+1)!}.
\tag{9}
\end{equation}
The first term is
\[
T_{j,1}=\frac{2e}{(2j+1)(2j+2)(2j+3)}.
\]
For \(\ell\geq 2\), every extra factorial factor is at least \(2j+1\), so
\[
0\leq j^3\sum_{\ell=2}^{\infty}
\ell\frac{(2j)!}{(2j+2\ell+1)!}
\leq j^3x_j^5
\left(\frac{2}{1-x_j^2}+
\frac{x_j^2}{(1-x_j^2)^2}\right),
\qquad x_j=\frac{1}{2j+1}.
\]
The right side is \(O(j^{-2})\).  Finally,
\[
j^3T_{j,1}\longrightarrow \frac{e}{4},
\]
which proves (2).

\section{Index and sign audit}

At the last backward step,
\[
P_3=\frac{123}{2},\qquad Q_3=396,\qquad R_3=36,
\qquad c_3=\frac1{20},
\]
so every required denominator is nonzero.  The upper terminal check from (7)
is
\[
\mu_{N-2}^{(N)}=\frac{1}{R_{N+1}},
\]
which makes the \(j=N+1\) equation recover \(\mu_{N+1}^{(N)}=1\) with the
prescribed zeros at \(N\) and \(N-1\).  For \(N=3\), (7) gives
\(\mu_1^{(3)}=1/480\) and \(\mu_0^{(3)}=11/480\), agreeing with direct
backward substitution.

All rearrangements above involve nonnegative terms.  No numerical evidence,
formal asymptotic ansatz, or unchecked symbolic computation is used in the
proof.

\begin{remark}
This theorem settles the \(c=1\) anchor only.  A closed form for the general
constant \(K(c)\), the source-term control in the broader recurrence program,
and the general coefficient-family classification remain open.
\end{remark}

\end{document}
